\documentclass[11pt]{article}

\usepackage{fullpage}
\usepackage{graphicx}
\usepackage{amsmath}
\usepackage{amssymb}
\usepackage{amsthm}
\usepackage{fancyvrb}
\usepackage{algorithm}
\usepackage{wrapfig}
\usepackage[noend]{algpseudocode}
\usepackage{xcolor}
\usepackage{tcolorbox}
\usepackage{hyperref}
\usepackage{caption}
\usepackage{subcaption}
\usepackage{booktabs}
\usepackage{microtype}


\tcbuselibrary{theorems, skins, breakable}

\theoremstyle{definition}
\newtheorem{definition}{Definition}
\theoremstyle{plain}
\newtheorem{theorem}{Theorem}
\newtheorem{lemma}{Lemma}
\newtheorem{corollary}{Corollary}
\newtheorem{remark}{Remark}
\newtheorem{example}{Example}
\newtheorem{proposition}{Proposition}

\parindent0in
\pagestyle{plain}
\thispagestyle{plain}
\graphicspath{{./}}

%% UPDATE MACRO DEFINITIONS %%
\newcommand{\myname}{Scribed By: Shreyas Mehta (2023101059)}
\newcommand{\mynami}{Harshit Lalwani (2023111028)}
\newcommand{\assignment}{Lecture \#12}
\newcommand{\duedate}{$23^{\text{rd}}$ February 2025}
\newtheorem*{thmtype}{Outline}
\newcommand{\statement}{%
This lecture covers: (1) convergence analysis of gradient descent for quadratic functions 
via a contraction recurrence; (2) the Kantorovich inequality and its role in establishing 
a linear convergence bound; (3) the condition number, its effect on convergence speed, 
and its practical relevance in machine learning; (4) the Rosenbrock function as a canonical 
ill-conditioned benchmark; (5) $L$-smooth functions, their characterisation via the Hessian, 
and the Descent Lemma as a stepping stone for convergence analysis of general smooth objectives.
}

% Convenience macros
\newcommand{\R}{\mathbb{R}}
\newcommand{\norm}[1]{\left\|#1\right\|}
\newcommand{\inner}[2]{\langle #1,\, #2\rangle}
\newcommand{\grad}{\nabla}

\begin{document}

\textbf{IIIT Hyderabad}\hfill\textbf{\myname}\\[0.01in]
\textbf{---}\hfill\textbf{\mynami}\\[0.01in]
\textbf{Optimization Methods (CS1.404)}\hfill\textbf{\assignment}\\[0.01in]
\textbf{Instructor: Dr.\ Naresh Manwani}\hfill\textbf{\duedate}\\
\smallskip\hrule\bigskip

\begin{thmtype}
\statement
\end{thmtype}

%%% ----------------------------------------------------------------
\section{Motivation: Comparing Convergence Across Step-size Strategies}
%%% ----------------------------------------------------------------

\begin{example}[Gradient Descent on $f(x,y)=x^2+2y^2$]\label{ex:gd_quad1}
Consider $f(x,y)=x^2+2y^2$ with optimal solution $(0,0)$ and optimal value $0$.
Starting from $(x_0,y_0)=(2,1)$, tolerance $\varepsilon=10^{-5}$:
\begin{itemize}
  \item \textbf{Exact line search:} Converges in \textbf{13 iterations} to 
    $(x^*,y^*)\approx(1.254\times10^{-6},\,-6.27\times10^{-6})$.
  \item \textbf{Constant step size} $t_k=0.1$: Converges in \textbf{58 iterations}.
    The step was too small, slowing convergence.
  \item \textbf{Backtracking line search} ($\tau=0.5,\,s=2,\,c_1=0.25$): Converges in
    \textbf{2 iterations} to the exact optimum.
    Inexact line search here outperforms exact line search.
\end{itemize}
\end{example}

\begin{example}[Gradient Descent on $f(x,y)=x^2+\tfrac{1}{100}y^2$]\label{ex:gd_quad2}
Consider $f(x,y)=x^2+\tfrac{1}{100}y^2$ with optimal solution $(0,0)$.
Starting from $(x_0,y_0)=(\tfrac{1}{100},1)$, $\varepsilon=10^{-5}$,
same backtracking parameters. The algorithm requires \textbf{201 iterations}.
\end{example}

Both functions are quadratics with minimum at the origin,
yet Example~\ref{ex:gd_quad1} converges in 2 steps while Example~\ref{ex:gd_quad2}
needs 201. The structural difference lies in the \emph{elongation} of the level sets:
$x^2+\tfrac{1}{100}y^2$ has highly elliptical contours, forcing gradient descent
to zigzag. The formal measure of this elongation is the \textbf{condition number}.

\medskip
\noindent\textbf{Backtracking line search --- geometric view.}
Figure~\ref{fig:boyd_backtrack} (reproduced from Boyd \& Vandenberghe~\cite{boyd},
Figure~9.1) illustrates the mechanism of the backtracking line search geometrically.
The curve shows $f$ restricted to the search ray $\{x + t\,\Delta x \mid t\geq 0\}$.
The lower dashed line is the linear extrapolation $f(x)+t\nabla f(x)^\top\Delta x$,
and the upper dashed line has slope scaled by the factor $\alpha\in(0,\tfrac{1}{2})$.
Backtracking accepts the first $t$ (starting from $t=1$ and multiplying by $\beta<1$
until acceptance) for which $f$ lies \emph{below} the upper dashed line, i.e.\
\[
  f(x+t\,\Delta x) \leq f(x) + \alpha\,t\,\nabla f(x)^\top\Delta x.
\]
The shaded region $[0, t_0]$ marks all acceptable step sizes; backtracking
terminates with $t=1$ (if already acceptable) or $t\in(\beta t_0, t_0]$.

\begin{figure}[h]
  \centering
  \includegraphics[width=0.55\textwidth]{boyd_fig9_1_backtracking.png}
  \caption{Backtracking line search (Boyd \& Vandenberghe~\cite{boyd}, Fig.~9.1).
    The curve is $f$ along the search ray. The lower dashed line is the full
    linear extrapolation; the upper dashed line has slope reduced by $\alpha$.
    The backtracking condition requires $f$ to lie below the upper dashed line,
    i.e.\ the step $t$ must lie in $[0,t_0]$.}
  \label{fig:boyd_backtrack}
\end{figure}

%%% ----------------------------------------------------------------
\section{Convergence of Steepest Descent with Exact Line Search for Quadratic Functions}
%%% ----------------------------------------------------------------

We now derive a precise bound for gradient descent on
\begin{equation}\label{eq:quad_prob}
  \min_{x\in\R^n}\; f(x) = x^\top A x,
\end{equation}
where $A\succ 0$ (symmetric positive definite). The unique minimiser is $x^*=0$
with $f^*=0$.

\subsection{Step Size from Exact Line Search}

The steepest descent direction is
\[
  d_k = -\grad f(x_k) = -2Ax_k.
\]
The exact line-search step size minimises $\phi(t)=f(x_k+td_k)$ over $t\geq 0$.
Expanding $\phi(t)$:
\begin{align*}
  \phi(t) &= (x_k+td_k)^\top A(x_k+td_k)\\
          &= x_k^\top Ax_k + 2t\,d_k^\top Ax_k + t^2\,d_k^\top Ad_k.
\end{align*}
Differentiating with respect to $t$ and setting to zero:
\[
  \phi'(t) = 2\,d_k^\top Ax_k + 2t\,d_k^\top Ad_k = 0
  \;\implies\;
  t_k = -\frac{d_k^\top Ax_k}{d_k^\top Ad_k}.
\]
Now substitute $d_k = -2Ax_k$:
\[
  d_k^\top Ax_k = (-2Ax_k)^\top Ax_k = -2x_k^\top A^\top Ax_k = -2x_k^\top A^2x_k.
\]
But we also have, using $d_k=-2Ax_k$ directly:
\[
  d_k^\top Ax_k = d_k^\top A\!\left(-\tfrac{1}{2}A^{-1}d_k\right) = -\tfrac{1}{2}d_k^\top d_k.
\]
(Here we used $x_k = -\tfrac{1}{2}A^{-1}d_k$, which follows from $d_k=-2Ax_k$.)
Substituting:
\begin{equation}\label{eq:tk}
  t_k = -\frac{-\tfrac{1}{2}d_k^\top d_k}{d_k^\top Ad_k}
      = \frac{d_k^\top d_k}{2\,d_k^\top Ad_k}.
\end{equation}

\subsection{Contraction Recurrence}

We compute $f(x_{k+1}) = f(x_k + t_k d_k)$ step by step.

\medskip\noindent\textbf{Step 1: Expand $f(x_k+t_kd_k)$.}
\begin{align}
  f(x_{k+1})
  &= (x_k+t_kd_k)^\top A(x_k+t_kd_k) \notag\\
  &= x_k^\top Ax_k + 2t_k\,d_k^\top Ax_k + t_k^2\,d_k^\top Ad_k. \label{eq:fk_step1}
\end{align}

\medskip\noindent\textbf{Step 2: Substitute $d_k^\top Ax_k = -\tfrac{1}{2}d_k^\top d_k$ and $t_k$ from~\eqref{eq:tk}.}
\begin{align}
  f(x_{k+1})
  &= x_k^\top Ax_k + 2\cdot\frac{d_k^\top d_k}{2\,d_k^\top Ad_k}\cdot\left(-\frac{d_k^\top d_k}{2}\right)
     + \left(\frac{d_k^\top d_k}{2\,d_k^\top Ad_k}\right)^{\!2} d_k^\top Ad_k \notag\\
  &= x_k^\top Ax_k - \frac{(d_k^\top d_k)^2}{2\,d_k^\top Ad_k}
     + \frac{(d_k^\top d_k)^2}{4\,d_k^\top Ad_k} \notag\\
  &= x_k^\top Ax_k - \frac{(d_k^\top d_k)^2}{4\,d_k^\top Ad_k}. \label{eq:fk_expand}
\end{align}

\medskip\noindent\textbf{Step 3: Rewrite as a contraction factor times $f(x_k)$.}

We use two identities. First, $f(x_k) = x_k^\top Ax_k$. Second,
since $d_k = -2Ax_k$, we have $x_k = -\tfrac{1}{2}A^{-1}d_k$, and therefore:
\[
  d_k^\top A^{-1}d_k
  = (-2Ax_k)^\top A^{-1}(-2Ax_k)
  = 4\,x_k^\top A^\top A^{-1} Ax_k
  = 4\,x_k^\top Ax_k
  = 4\,f(x_k).
\]
Thus $x_k^\top Ax_k = \tfrac{1}{4}d_k^\top A^{-1}d_k$. Substituting into~\eqref{eq:fk_expand}:
\begin{align}
  f(x_{k+1})
  &= x_k^\top Ax_k - \frac{(d_k^\top d_k)^2}{4\,d_k^\top Ad_k} \notag\\
  &= x_k^\top Ax_k\left(1 - \frac{(d_k^\top d_k)^2}{4\,d_k^\top Ad_k \cdot x_k^\top Ax_k}\right) \notag\\
  &= x_k^\top Ax_k\left(1 - \frac{(d_k^\top d_k)^2}{4\,d_k^\top Ad_k \cdot \tfrac{1}{4}d_k^\top A^{-1}d_k}\right) \notag\\
  &= \left(1 - \frac{(d_k^\top d_k)^2}{(d_k^\top Ad_k)(d_k^\top A^{-1}d_k)}\right)f(x_k).
  \label{eq:ratio_form}
\end{align}

\noindent\textbf{Intuition.}
The ratio $\rho_k = (d_k^\top d_k)^2/[(d_k^\top Ad_k)(d_k^\top A^{-1}d_k)]$ measures
how much progress is made each iteration. When $A=I$ (perfectly spherical contours),
$\rho_k=1$ and convergence is immediate. When $A$ is far from a multiple of $I$
(elongated contours), $\rho_k$ can be close to zero---very little progress per step.
The Kantorovich inequality provides a \emph{universal lower bound} on $\rho_k$.

%%% ----------------------------------------------------------------
\section{Kantorovich Inequality and Linear Convergence}
%%% ----------------------------------------------------------------

\begin{lemma}[Kantorovich Inequality]\label{lem:kantorovich}
Let $A$ be a positive definite $n\times n$ matrix with
$m=\lambda_{\min}(A)$ and $M=\lambda_{\max}(A)$. Then for any
$0\neq x\in\R^n$:
\begin{equation}\label{eq:kantorovich}
  \frac{(x^\top x)^2}{(x^\top Ax)(x^\top A^{-1}x)}
  \;\geq\; \frac{4Mm}{(M+m)^2}.
\end{equation}
\end{lemma}

\begin{proof}
Let $m = \lambda_{\min}(A)$ and $M = \lambda_{\max}(A)$.

$A$ is symmetric positive definite, the spectral theorem guarantees the existence
of an orthogonal matrix $U$ such that:
\[
  A = U\Lambda U^\top,
\]
where $\Lambda = \operatorname{diag}(\lambda_1, \lambda_2, \ldots, \lambda_n)$
is the diagonal matrix of eigenvalues, with $0 < m \leq \lambda_i \leq M$ for all $i$.

\[
  A^{-1} = (U\Lambda U^\top)^{-1} = (U^\top)^{-1}\Lambda^{-1}U^{-1} = U\Lambda^{-1}U^\top,
\]
where $\Lambda^{-1} = \operatorname{diag}(1/\lambda_1, 1/\lambda_2, \ldots, 1/\lambda_n)$.
Each entry $1/\lambda_i$ is well-defined since $\lambda_i > 0$.


Substituting the decompositions of $A$ and $A^{-1}$:
\begin{align*}
  A + Mm\,A^{-1}
  &= U\Lambda U^\top + Mm\cdot U\Lambda^{-1}U^\top \\
  &= U\bigl(\Lambda + Mm\,\Lambda^{-1}\bigr)U^\top,
\end{align*}
where the factorisation is valid because $U$ is the same orthogonal matrix for both terms.

The expression $U(\Lambda + Mm\,\Lambda^{-1})U^\top$ is the eigendecomposition of $A+Mm\,A^{-1}$.
Since $\Lambda + Mm\,\Lambda^{-1}$ is diagonal, its diagonal entries are the eigenvalues:
\[
  \mu_i = \lambda_i + \frac{Mm}{\lambda_i}, \qquad i = 1, 2, \ldots, n.
\]

We use the Rayleigh quotient characterisation:
\[
  \lambda_{\max}(A + Mm\,A^{-1}) = \sup_{x\neq 0}\frac{x^\top(A+Mm\,A^{-1})x}{x^\top x}.
\]
Since the eigenvalues of $A+Mm\,A^{-1}$ are $\{\mu_i\}$, and the supremum of the
Rayleigh quotient equals the largest eigenvalue, we need to find $\max_i \mu_i$.
Define $\varphi:[m,M]\to\R$ by $\varphi(t)=t+Mm/t$. Computing its derivatives:
\[
  \varphi'(t) = 1 - \frac{Mm}{t^2}, \qquad \varphi''(t) = \frac{2Mm}{t^3} > 0,
\]
so $\varphi$ is strictly convex on $[m,M]$, and its maximum over the closed interval
is attained at one of the endpoints. Computing both:
\[
  \varphi(m) = m + \frac{Mm}{m} = m + M, \qquad
  \varphi(M) = M + \frac{Mm}{M} = M + m.
\]
Both endpoints give $M+m$, so for all $i$:
\[
  \mu_i = \varphi(\lambda_i) \leq M + m,
\]
and hence:
\[
  \lambda_{\max}(A + Mm\,A^{-1}) = \max_i\,\mu_i \leq M + m.
\]

For any symmetric matrix $B$, the Rayleigh quotient satisfies:
\[
  \frac{x^\top Bx}{x^\top x} \leq \lambda_{\max}(B) \quad\forall\;x\neq 0.
\]
Applying this to $B = A + Mm\,A^{-1}$ with $\lambda_{\max}(B)\leq M+m$:
\[
  \frac{x^\top(A + Mm\,A^{-1})x}{x^\top x} \leq M + m.
\]
Multiplying both sides by $x^\top x = \norm{x}^2 > 0$:
\begin{equation}\label{eq:Bbound}
  x^\top(A + Mm\,A^{-1})x \leq (M+m)\norm{x}^2,
\end{equation}
which is equivalent to the positive semidefinite ordering $A + Mm\,A^{-1} \preceq (M+m)I$.


Expanding the left side of~\eqref{eq:Bbound}:
\begin{equation}\label{eq:quad_form_ineq}
  x^\top Ax + Mm\,(x^\top A^{-1}x) \leq (M+m)\,\norm{x}^2.
\end{equation}


For any two non-negative reals $\alpha,\beta\geq 0$, the AM-GM inequality states:
\[
  \sqrt{\alpha\beta} \leq \frac{\alpha+\beta}{2}
  \quad\Longleftrightarrow\quad
  \alpha\beta \leq \left(\frac{\alpha+\beta}{2}\right)^2.
\]
Set $\alpha = x^\top Ax \geq 0$ and $\beta = Mm\cdot x^\top A^{-1}x \geq 0$:
\begin{align}
  \bigl(x^\top Ax\bigr)\cdot Mm\,\bigl(x^\top A^{-1}x\bigr)
  &\leq \left(\frac{x^\top Ax + Mm\,x^\top A^{-1}x}{2}\right)^{\!2}. \label{eq:amgm_step}
\end{align}
Applying~\eqref{eq:quad_form_ineq} to bound the right side of~\eqref{eq:amgm_step}:
\begin{equation}\label{eq:amgm_applied}
  \bigl(x^\top Ax\bigr)\cdot Mm\,\bigl(x^\top A^{-1}x\bigr)
  \leq \left(\frac{(M+m)\norm{x}^2}{2}\right)^{\!2}
  = \frac{(M+m)^2}{4}\,\norm{x}^4.
\end{equation}

Dividing both sides of~\eqref{eq:amgm_applied} by $Mm\,\bigl(x^\top Ax\bigr)\bigl(x^\top A^{-1}x\bigr) > 0$:
\[
  1 \leq \frac{(M+m)^2\,\norm{x}^4}{4\,Mm\,\bigl(x^\top Ax\bigr)\bigl(x^\top A^{-1}x\bigr)},
\]
which rearranges to:
\[
  \frac{\norm{x}^4}{\bigl(x^\top Ax\bigr)\bigl(x^\top A^{-1}x\bigr)} \geq \frac{4Mm}{(M+m)^2}.
\]
Since $\norm{x}^4 = (x^\top x)^2$, this is exactly the claimed inequality~\eqref{eq:kantorovich}. \qedhere
\end{proof}

\noindent\textbf{Intuition.}
The Kantorovich inequality says: no matter which direction $x$ we are in,
the alignment ratio $\rho=(x^\top x)^2/[(x^\top Ax)(x^\top A^{-1}x)]$
is always at least $4Mm/(M+m)^2$. The tighter the eigenvalue spread
(i.e.\ smaller $M/m$), the larger this lower bound, meaning more progress
is guaranteed per step.

\begin{lemma}[Linear Convergence for Quadratic Minimisation]\label{lem:linear_conv}
Let $\{x_k\}_{k\geq 0}$ be the sequence generated by gradient descent with
exact line search for problem~\eqref{eq:quad_prob}. Then for any $k=0,1,2,\ldots$,
\begin{equation}\label{eq:linear_conv}
  f(x_{k+1})
  \;\leq\; \left(\frac{M-m}{M+m}\right)^{\!2} f(x_k),
\end{equation}
where $M=\lambda_{\max}(A)$ and $m=\lambda_{\min}(A)$.
\end{lemma}

\begin{proof}

Substituting $x = d_k$ into~\eqref{eq:kantorovich}:
\[
  \frac{(d_k^\top d_k)^2}{(d_k^\top Ad_k)(d_k^\top A^{-1}d_k)} \geq \frac{4Mm}{(M+m)^2}.
\]

Since $0 \leq \frac{(d_k^\top d_k)^2}{(d_k^\top Ad_k)(d_k^\top A^{-1}d_k)} \leq 1$,
the factor $\bigl(1 - \rho_k\bigr)$ in~\eqref{eq:ratio_form} satisfies:
\[
  1 - \frac{(d_k^\top d_k)^2}{(d_k^\top Ad_k)(d_k^\top A^{-1}d_k)}
  \leq 1 - \frac{4Mm}{(M+m)^2}.
\]

\begin{align}
  1 - \frac{4Mm}{(M+m)^2}
  &= \frac{(M+m)^2 - 4Mm}{(M+m)^2} \notag\\
  &= \frac{M^2 + 2Mm + m^2 - 4Mm}{(M+m)^2} \notag\\
  &= \frac{M^2 - 2Mm + m^2}{(M+m)^2} \notag\\
  &= \frac{(M-m)^2}{(M+m)^2}
   = \left(\frac{M-m}{M+m}\right)^{\!2}. \label{eq:c_derivation}
\end{align}


Substituting into~\eqref{eq:ratio_form}:
\[
  f(x_{k+1})
  = \left(1 - \frac{(d_k^\top d_k)^2}{(d_k^\top Ad_k)(d_k^\top A^{-1}d_k)}\right)f(x_k)
  \leq \left(\frac{M-m}{M+m}\right)^{\!2}f(x_k). \qedhere
\]
\end{proof}

\begin{corollary}[Geometric Decay]\label{cor:geometric}
Applying Lemma~\ref{lem:linear_conv} repeatedly:
\[
  f(x_k) \leq c^k\,f(x_0), \qquad c = \left(\frac{M-m}{M+m}\right)^{\!2}
  = \left(\frac{\kappa-1}{\kappa+1}\right)^{\!2},
\]
where $\kappa=M/m$. The sequence $\{f(x_k)\}$ is bounded above by a geometrically
decreasing sequence; this is called \textbf{linear} (or $Q$-linear) convergence.
\end{corollary}

\noindent\textbf{Summary table.} Since $c$ is increasing in $\kappa$:

\begin{center}
\begin{tabular}{cccc}
\toprule
$\kappa$ & $c=\left(\frac{\kappa-1}{\kappa+1}\right)^2$ & Iterations to halve $f$ & Character\\
\midrule
$1$    & $0$       & $1$     & Perfectly conditioned \\
$2$    & $\approx 0.11$ & $\approx 1$  & Well-conditioned \\
$10$   & $\approx 0.67$ & $\approx 2$  & Mild ill-conditioning \\
$100$  & $\approx 0.960$ & $\approx 17$ & Ill-conditioned \\
$1000$ & $\approx 0.996$ & $\approx 173$ & Severely ill-conditioned \\
\bottomrule
\end{tabular}
\end{center}

%%% ----------------------------------------------------------------
\section{Condition Number}
%%% ----------------------------------------------------------------

\begin{definition}[Condition Number]\label{def:cond}
Let $A$ be an $n\times n$ positive definite matrix. Its \textbf{condition number} is
\[
  \kappa(A) = \frac{\lambda_{\max}(A)}{\lambda_{\min}(A)}.
\]
Matrices with large $\kappa$ are called \emph{ill-conditioned};
those with $\kappa$ close to $1$ are \emph{well-conditioned}.
\end{definition}

\noindent\textbf{Geometric interpretation.}
$\kappa(A)$ equals the ratio of the longest to shortest semi-axis of
the ellipsoid $\{x : x^\top Ax \leq 1\}$. A large $\kappa$ means highly
elongated elliptical contours. The gradient at a point on an elongated
ellipsoid points nearly perpendicular to the direction of the minimum,
so steepest descent zigzags inefficiently.

\medskip
\noindent\textit{Note.} Additional discussion on sublevel-set condition numbers,
convergence speed, and feature normalisation is moved to
Annexure~\ref{app:cond_annexure}.

%%% ----------------------------------------------------------------
\section{The Rosenbrock Function: A Canonical Ill-Conditioned Benchmark}
%%% ----------------------------------------------------------------

\begin{example}[Rosenbrock Function]\label{ex:rosenbrock}
The \textbf{Rosenbrock function} is:
\[
  f(x_1,x_2) = 100(x_2-x_1^2)^2 + (1-x_1)^2.
\]
\begin{itemize}
  \item \textbf{Optimum:} $(x_1^*,x_2^*)=(1,1)$, $f^*=0$.
  \item \textbf{Gradient:}
    $\grad f(x)=\bigl(-400x_1(x_2-x_1^2)-2(1-x_1),\;200(x_2-x_1^2)\bigr)^\top$.
  \item \textbf{Hessian:}
    \[
      \grad^2 f(x) = \begin{pmatrix}-400x_2+1200x_1^2+2 & -400x_1\\-400x_1 & 200\end{pmatrix}.
    \]
  \item \textbf{Hessian at optimum:}
    $\grad^2 f(1,1)=\begin{pmatrix}802&-400\\-400&200\end{pmatrix}$, $\quad\kappa\approx 2508$.
  \item Backtracking gradient descent from $x_0=(2,5)^\top$ requires \textbf{6890 iterations}.
\end{itemize}
The \emph{banana-shaped} contour lines (Figure~\ref{fig:rosenbrock}) surround the
unique minimum and force the gradient to point almost perpendicularly to the valley,
causing persistent zigzagging.
\end{example}

\begin{figure}[h]
  \centering
  \includegraphics[width=0.60\textwidth]{rosenbrock_contour.pdf}
  \caption{Contour lines of the Rosenbrock function and the GD iterates
    (red path) from $x_0=(2,5)^\top$. The banana-shaped valley makes steepest
    descent take thousands of tiny zigzag steps before reaching $(1,1)$ (black star).}
  \label{fig:rosenbrock}
\end{figure}

\noindent\textbf{Why is Rosenbrock a standard benchmark?}
Any new optimisation algorithm is routinely tested on the Rosenbrock function
(and its $n$-dimensional generalisation
$\sum_{i=1}^{n-1}\bigl[100(x_{i+1}-x_i^2)^2+(1-x_i)^2\bigr]$) precisely
because its severe ill-conditioning ($\kappa\approx 2508$ in 2D) exposes
weaknesses in gradient-based methods. An algorithm that converges quickly
on Rosenbrock is considered robust.

\begin{example}[A Related Ill-Conditioned Function: Booth's Function \textmd{\textcolor{teal}{[Extra Example]}}]\label{ex:booth}
Another test function useful for comparison is \textbf{Booth's function}:
\[
  f(x_1,x_2) = (x_1+2x_2-7)^2 + (2x_1+x_2-5)^2.
\]
The minimum is at $(1,3)$ with $f^*=0$.
The Hessian is constant: $\grad^2 f = \begin{pmatrix}10&8\\8&10\end{pmatrix}$,
with eigenvalues $2$ and $18$, so $\kappa=9$. Gradient descent converges much
faster here than on Rosenbrock, illustrating again the effect of $\kappa$.
\end{example}

%%% ----------------------------------------------------------------
\section{$L$-Smooth Functions}
%%% ----------------------------------------------------------------

Everything up to this point assumed a \emph{quadratic} objective where the
Hessian is \emph{constant}. For general differentiable functions, we need
a condition that bounds the Hessian \emph{everywhere}, not just at one point.
This leads to the notion of $L$-smoothness.

\subsection{Definition and Interpretation}

\begin{definition}[$L$-Smooth Function]\label{def:L_smooth}
A continuously differentiable function $f:\R^n\to\R$ is $L$-\textbf{smooth}
(or has an \emph{$L$-Lipschitz gradient}) if there exists $L>0$ such that
\[
  \norm{\grad f(x) - \grad f(y)} \leq L\norm{x-y}
  \quad\forall\; x,y\in\R^n.
\]
The class of such functions is denoted $\mathcal{C}^{1,1}_L(\R^n)$.
\end{definition}

\noindent\textbf{Intuition (from lecture).}
The gradient of a function tells us the direction of steepest ascent.
If the gradient could change arbitrarily quickly, knowing $\grad f(x)$
would give us no information about $f$ near $x$, making gradient descent
unreliable. $L$-smoothness guarantees that the gradient is \emph{stable}:
moving $\delta$ in input space can change the gradient by at most $L\delta$.
This means gradient information is trustworthy within a neighbourhood,
and we can safely decrease $f$ by stepping against the gradient.

\begin{remark}[Infinitely many valid $L$]
If $\norm{\grad f(x)-\grad f(y)}\leq L\norm{x-y}$ holds for some $L$,
it also holds for any $L'>L$. So there are infinitely many valid Lipschitz
constants; we are typically interested in the \emph{smallest} one.
\end{remark}

\begin{remark}[Slide Correction --- Slide 17]\label{rem:slide_correction}
In the slides, the quadratic example writes
\[
  \|f(x)-f(y)\| = 2\|A(x-y)\| \leq 2\|A\|\cdot\|x-y\|.
\]
\textbf{This is a typographical error.} The correct statement involves
the \emph{gradient} difference, not the function value difference:
\[
  \|\grad f(x)-\grad f(y)\| = 2\|A(x-y)\| \leq 2\|A\|\cdot\|x-y\|.
\]
Hence the Lipschitz constant of $\grad f$ for the quadratic is $L=2\|A\|$.
Note that $f(x)-f(y)$ is a scalar so its ``norm'' would just be an absolute value,
but more importantly the Lipschitz condition for $L$-smoothness applies to
\emph{gradients}, not function values.
\end{remark}

\subsection{Examples of $L$-Smooth Functions}

\begin{example}[Linear Functions]\label{ex:linear_smooth}
Let $f(x)=a^\top x$ for some $a\in\R^n$.
Then $\grad f(x)=a$ for all $x$, so
$\norm{\grad f(x)-\grad f(y)}=0\leq 0\cdot\norm{x-y}$.
Hence $f\in\mathcal{C}^{1,1}_0$. The gradient is constant and no smoothness constant is needed.
\end{example}

\begin{example}[Quadratic Functions]\label{ex:quad_smooth}
Let $f(x)=x^\top Ax+2b^\top x+c$ with $A$ symmetric. Then $\grad f(x)=2Ax+2b$, so
\[
  \norm{\grad f(x)-\grad f(y)} = 2\norm{A(x-y)} \leq 2\norm{A}\norm{x-y}.
\]
Hence $f\in\mathcal{C}^{1,1}_{2\|A\|}$. For a diagonal matrix
$A=\operatorname{diag}(\lambda_1,\ldots,\lambda_n)$, $\norm{A}=\lambda_{\max}(A)$,
so $L=2\lambda_{\max}(A)$.
\end{example}

\begin{example}[$f(x)=\sqrt{1+x^2}$]\label{ex:sqrt_smooth}
We have $f'(x)=x/\sqrt{1+x^2}$ and
\[
  f''(x) = \frac{1}{(1+x^2)^{3/2}}.
\]
Since $0\leq f''(x)\leq 1$ for all $x\in\R$, Theorem~\ref{thm:hessian_equiv} below
gives $f\in\mathcal{C}^{1,1}_1$.
\end{example}

\begin{example}[Logistic Loss \textmd{\textcolor{teal}{[Extra Example]}}]\label{ex:logistic_smooth}
The logistic loss used in binary classification is
\[
  f(w) = \frac{1}{n}\sum_{i=1}^n \log\!\left(1+e^{-y_i x_i^\top w}\right),
\]
where $y_i\in\{-1,+1\}$. This is $L$-smooth because the sigmoid
$\sigma(t)=1/(1+e^{-t})$ satisfies $0\leq\sigma''(t)=\sigma(t)(1-\sigma(t))\leq 1/4$.
By the chain rule, the Hessian of $f$ satisfies
\[
  \norm{\grad^2 f(w)} \leq \frac{1}{4n}\sum_{i=1}^n\norm{x_i}^2.
\]
Hence $f\in\mathcal{C}^{1,1}_L$ with $L=\frac{1}{4n}\sum_{i=1}^n\norm{x_i}^2$.
This means gradient descent with step size $t=1/L$ is guaranteed to converge.
\end{example}

\begin{example}[$f(x)=e^{x}$ is \emph{not} $L$-smooth on $\R$ \textmd{\textcolor{teal}{[Extra Example]}}]\label{ex:exp_not_smooth}
Consider $f(x)=e^x$ on all of $\R$. Then $f''(x)=e^x$, which is unbounded
as $x\to\infty$. By Theorem~\ref{thm:hessian_equiv}, $f\notin\mathcal{C}^{1,1}_L$ for any finite $L$.
Note, however, that $e^x$ \emph{is} $L$-smooth on any bounded interval $[-M,M]$
with $L=e^M$. This shows that $L$-smoothness is a \emph{global} condition
when stated over all of $\R^n$.
\end{example}

\subsection{$L$-Smoothness vs.\ Lipschitz Continuity of $f$}

A common confusion: is $L$-smoothness (Lipschitz \emph{gradient})
the same as Lipschitz continuity of $f$ itself?

\begin{proposition}[$L$-Smoothness is Strictly Stronger]\label{prop:smooth_implies_lip}
$L$-smoothness is a \textbf{strictly stronger} condition than
Lipschitz continuity of $f$. Specifically:
\begin{enumerate}
  \item[(a)] If $f\in\mathcal{C}^{1,1}_L$ and $\norm{\grad f}$ is bounded
    by some $G$, then $f$ is Lipschitz continuous with constant $G$.
  \item[(b)] Lipschitz continuity of $f$ does \emph{not} imply $L$-smoothness
    (even differentiability is not guaranteed).
\end{enumerate}
\end{proposition}

\begin{proof}
\textbf{(a): $L$-smooth + bounded gradient $\Rightarrow$ Lipschitz.}

Let $f\in\mathcal{C}^{1,1}_L$ and suppose $\norm{\grad f(z)}\leq G$ for all $z\in\R^n$.
By the fundamental theorem of calculus along the line segment from $x$ to $y$:
\[
  f(y) - f(x) = \int_0^1 \inner{\grad f\bigl(x+t(y-x)\bigr)}{y-x}\,dt.
\]
Taking absolute values and applying Cauchy-Schwarz:
\begin{align*}
  |f(y)-f(x)|
  &= \left|\int_0^1 \inner{\grad f\bigl(x+t(y-x)\bigr)}{y-x}\,dt\right| \\
  &\leq \int_0^1 \norm{\grad f\bigl(x+t(y-x)\bigr)}\cdot\norm{y-x}\,dt \\
  &\leq \int_0^1 G\,\norm{y-x}\,dt
  = G\,\norm{y-x}.
\end{align*}
Hence $f$ is Lipschitz with constant $G$.

\medskip
\textbf{(b) Counter-example: Lipschitz $\not\Rightarrow$ $L$-smooth.}

Let $f:\R\to\R$ be $f(x)=|x|$. We verify it is Lipschitz:
\[
  |f(x)-f(y)| = \bigl||x|-|y|\bigr| \leq |x-y|
  \quad\text{(reverse triangle inequality)},
\]
so $f$ is Lipschitz with constant $1$. However, $f$ is \emph{not differentiable} at $x=0$:
\[
  \lim_{h\to 0^+}\frac{f(0+h)-f(0)}{h} = \lim_{h\to 0^+}\frac{h}{h} = 1 \neq -1 = \lim_{h\to 0^-}\frac{-h}{h}.
\]
Since $L$-smoothness requires $f$ to be continuously differentiable with a Lipschitz gradient,
and $f(x)=|x|$ is not even differentiable at $x=0$, it cannot be $L$-smooth for any $L\geq 0$.
\end{proof}

Figure~\ref{fig:lipschitz_smooth} illustrates the distinction: $|x|$ is Lipschitz
but its subgradient jumps discontinuously at the origin,
while $\sqrt{x^2+\varepsilon}$ is a smooth approximation with a Lipschitz gradient.

\begin{figure}[h]
  \centering
  \includegraphics[width=0.95\textwidth]{lipschitz_vs_smooth.pdf}
  \caption{Left: $f(x)=|x|$ (Lipschitz but not smooth, blue) vs.\
    $f(x)=\sqrt{x^2+0.1}$ ($L$-smooth, red). Right: the subgradient of $|x|$
    has a jump discontinuity at the origin (violating Lipschitz continuity
    of the gradient), whereas $\nabla\sqrt{x^2+0.1}$ is continuous and bounded.
    This illustrates why $L$-smoothness is a strictly stronger condition.}
  \label{fig:lipschitz_smooth}
\end{figure}

\subsection{Equivalent Characterisation via the Hessian}

\begin{theorem}[Beck, Theorem~4.20]\label{thm:hessian_equiv}
Let $f$ be twice continuously differentiable over $\R^n$. The following
are equivalent:
\begin{enumerate}
  \item[(a)] $f\in\mathcal{C}^{1,1}_L(\R^n)$.
  \item[(b)] $\norm{\grad^2 f(x)}\leq L$ for all $x\in\R^n$.
\end{enumerate}
\end{theorem}

\begin{proof}
\textbf{(b)$\Rightarrow$(a): Bounded Hessian implies $L$-Lipschitz gradient.}

Suppose $\norm{\grad^2 f(x)}\leq L$ for all $x\in\R^n$.
Fix any $x,y\in\R^n$. By the fundamental theorem of calculus applied
to $g(t)=\grad f(x+t(y-x))$:
\[
  \grad f(y) - \grad f(x)
  = \int_0^1 \frac{d}{dt}\grad f(x+t(y-x))\,dt
  = \int_0^1 \grad^2 f\bigl(x+t(y-x)\bigr)(y-x)\,dt.
\]
Taking norms and using sub-multiplicativity of the operator norm ($\norm{Bv}\leq\norm{B}\norm{v}$):
\begin{align*}
  \norm{\grad f(y)-\grad f(x)}
  &= \left\|\int_0^1 \grad^2 f\bigl(x+t(y-x)\bigr)(y-x)\,dt\right\|\\
  &\leq \int_0^1 \left\|\grad^2 f\bigl(x+t(y-x)\bigr)(y-x)\right\|\,dt
  \quad\text{(triangle inequality for integrals)}\\
  &\leq \int_0^1 \left\|\grad^2 f\bigl(x+t(y-x)\bigr)\right\|\cdot\norm{y-x}\,dt
  \quad\text{(operator norm bound)}\\
  &\leq \int_0^1 L\,\norm{y-x}\,dt
  \quad\text{(using assumption (b))}\\
  &= L\,\norm{y-x}.
\end{align*}
Hence $f\in\mathcal{C}^{1,1}_L(\R^n)$.

\medskip
\textbf{(a)$\Rightarrow$(b): $L$-Lipschitz gradient implies bounded Hessian.}

Suppose $f\in\mathcal{C}^{1,1}_L(\R^n)$. Fix $x\in\R^n$,
and let $d\in\R^n$ be an arbitrary direction with $\norm{d}=1$. For any $\alpha>0$:
\[
  \grad f(x+\alpha d)-\grad f(x)
  = \int_0^1 \frac{d}{dt}\grad f(x+t\alpha d)\,dt
  = \int_0^1 \grad^2 f(x+t\alpha d)\cdot(\alpha d)\,dt
  = \alpha\int_0^1 \grad^2 f(x+t\alpha d)\,d\,dt.
\]
Taking norms on both sides:
\[
  \left\|\int_0^1 \grad^2 f(x+t\alpha d)\,d\,dt\right\|
  = \frac{1}{\alpha}\norm{\grad f(x+\alpha d)-\grad f(x)}
  \leq \frac{1}{\alpha}\cdot L\norm{\alpha d}
  = L\norm{d} = L.
\]
Dividing by $\norm{d}=1$ and taking the limit $\alpha\to 0^+$,
by continuity of $\grad^2 f$:
\[
  \norm{\grad^2 f(x)\,d}
  = \lim_{\alpha\to 0^+}\left\|\int_0^1 \grad^2 f(x+t\alpha d)\,d\,dt\right\|
  \leq L.
\]
Since $d$ was an arbitrary unit vector, we conclude $\norm{\grad^2 f(x)} = \sup_{\norm{d}=1}\norm{\grad^2 f(x)d} \leq L$.
\end{proof}

\noindent\textbf{Intuition.} This theorem says $L$-smoothness is \emph{exactly}
the condition that the curvature (Hessian) is bounded everywhere.
When $f$ is quadratic, the Hessian is constant and equals $2A$;
hence $L=2\norm{A}=2\lambda_{\max}(A)$, connecting directly back to the quadratic analysis.
For non-quadratic functions, $\norm{\grad^2 f(x)}$ varies with $x$,
so we need it to be bounded \emph{uniformly} over all $x$.

%%% ----------------------------------------------------------------
\section{The Descent Lemma}
%%% ----------------------------------------------------------------

\begin{lemma}[Descent Lemma, Beck Lemma~4.22]\label{lem:descent}
Let $f\in\mathcal{C}^{1,1}_L(\R^n)$. Then for any $x,y\in\R^n$:
\begin{equation}\label{eq:descent_upper}
  f(y) \leq f(x) + \grad f(x)^\top(y-x) + \frac{L}{2}\norm{x-y}^2.
\end{equation}
Moreover, the proof establishes the two-sided sandwich:
\begin{equation}\label{eq:descent_sandwich}
  f(x) + \grad f(x)^\top(y-x) - \frac{L}{2}\norm{x-y}^2
  \;\leq\; f(y)
  \;\leq\; f(x) + \grad f(x)^\top(y-x) + \frac{L}{2}\norm{x-y}^2.
\end{equation}
\end{lemma}

\begin{proof}
\textbf{Step 1: Express $f(y)-f(x)$ via the fundamental theorem of calculus.}

Let $\gamma(t) = x + t(y-x)$ for $t\in[0,1]$. Then:
\[
  f(y) - f(x)
  = \int_0^1 \frac{d}{dt}f(\gamma(t))\,dt
  = \int_0^1 \inner{\grad f\bigl(x+t(y-x)\bigr)}{y-x}\,dt.
\]

\medskip\noindent\textbf{Step 2: Isolate the deviation from the linear approximation.}

Add and subtract $\inner{\grad f(x)}{y-x}$ inside the integral:
\begin{align*}
  f(y) - f(x) - \inner{\grad f(x)}{y-x}
  &= \int_0^1 \inner{\grad f\bigl(x+t(y-x)\bigr)}{y-x}\,dt
     - \inner{\grad f(x)}{y-x}\\
  &= \int_0^1 \inner{\grad f\bigl(x+t(y-x)\bigr) - \grad f(x)}{y-x}\,dt.
\end{align*}

\medskip\noindent\textbf{Step 3: Bound the deviation using Cauchy-Schwarz and $L$-smoothness.}

Taking the absolute value:
\begin{align*}
  \bigl|f(y)-f(x)-\inner{\grad f(x)}{y-x}\bigr|
  &= \left|\int_0^1 \inner{\grad f\bigl(x+t(y-x)\bigr)-\grad f(x)}{y-x}\,dt\right|\\
  &\leq \int_0^1 \left|\inner{\grad f\bigl(x+t(y-x)\bigr)-\grad f(x)}{y-x}\right|\,dt
  \quad\text{(triangle ineq.)}\\
  &\leq \int_0^1 \norm{\grad f\bigl(x+t(y-x)\bigr)-\grad f(x)}\cdot\norm{y-x}\,dt
  \quad\text{(Cauchy-Schwarz)}\\
  &\leq \int_0^1 L\,\norm{t(y-x)}\cdot\norm{y-x}\,dt
  \quad\text{($L$-smoothness: }\norm{\grad f(x+h)-\grad f(x)}\leq L\norm{h}\text{)}\\
  &= \int_0^1 L\,t\,\norm{y-x}^2\,dt
  \quad\text{(since }\norm{t(y-x)}=t\norm{y-x}\text{)}\\
  &= L\norm{y-x}^2\int_0^1 t\,dt\\
  &= \frac{L}{2}\norm{y-x}^2.
\end{align*}

\medskip\noindent\textbf{Step 4: Obtain the two-sided bound.}

The inequality $|a|\leq b$ is equivalent to $-b\leq a\leq b$.
Applying this with $a = f(y)-f(x)-\inner{\grad f(x)}{y-x}$ and $b=\frac{L}{2}\norm{y-x}^2$:
\[
  -\frac{L}{2}\norm{y-x}^2
  \leq f(y)-f(x)-\inner{\grad f(x)}{y-x}
  \leq \frac{L}{2}\norm{y-x}^2.
\]
Adding $f(x)+\inner{\grad f(x)}{y-x}$ to all three sides gives~\eqref{eq:descent_sandwich},
and the upper bound is~\eqref{eq:descent_upper}.
\end{proof}

\noindent\textbf{Intuition.}
The Descent Lemma says that an $L$-smooth function is \emph{sandwiched}
between a linearisation minus a quadratic correction and
a linearisation plus a quadratic correction (Figure~\ref{fig:descent_lemma}).
The upper bound~\eqref{eq:descent_upper} is the workhorse: it says
``the function cannot curve upward faster than $L/2$ times the squared distance.''
This is the $L$-smooth analogue of the constant-Hessian quadratic upper bound
used in the quadratic analysis, and it is the key ingredient in proving
convergence of gradient descent for general smooth objectives.

\begin{figure}[h]
  \centering
  \includegraphics[width=0.60\textwidth]{descent_lemma.pdf}
  \caption{Descent Lemma for $f(x)=\sin(x)+\tfrac{1}{2}x^2$ (blue, $L=3$).
    The upper bound $Q^+$ (red) and lower bound $Q^-$ (green dashed) are
    quadratics centred at the expansion point $x$ (black dot).
    The function lies entirely within the sandwich for all $y$.}
  \label{fig:descent_lemma}
\end{figure}

\subsection{Immediate Consequence: Gradient Step Decreases $f$}

Setting $y = x - \frac{1}{L}\grad f(x)$ in the upper bound~\eqref{eq:descent_upper}
and computing each term:
\begin{align*}
  y - x &= -\frac{1}{L}\grad f(x), \\
  \grad f(x)^\top(y-x) &= \grad f(x)^\top\!\left(-\frac{1}{L}\grad f(x)\right) = -\frac{1}{L}\norm{\grad f(x)}^2, \\
  \norm{x-y}^2 &= \norm{\tfrac{1}{L}\grad f(x)}^2 = \frac{1}{L^2}\norm{\grad f(x)}^2.
\end{align*}
Substituting into~\eqref{eq:descent_upper}:
\begin{align*}
  f\!\left(x - \tfrac{1}{L}\grad f(x)\right)
  &\leq f(x) + \left(-\frac{1}{L}\norm{\grad f(x)}^2\right)
     + \frac{L}{2}\cdot\frac{1}{L^2}\norm{\grad f(x)}^2\\
  &= f(x) - \frac{1}{L}\norm{\grad f(x)}^2 + \frac{1}{2L}\norm{\grad f(x)}^2\\
  &= f(x) - \frac{2}{2L}\norm{\grad f(x)}^2 + \frac{1}{2L}\norm{\grad f(x)}^2\\
  &= f(x) - \frac{1}{2L}\norm{\grad f(x)}^2.
\end{align*}
This is the \textbf{Sufficient Decrease} property: a gradient step of size $1/L$
always reduces the objective by at least $\frac{1}{2L}\norm{\grad f(x)}^2$.
The formal lemma for all valid step-size strategies (constant $1/L$, exact line
search, backtracking) was \emph{stated but not proved} in this lecture; the proof
is deferred to the next class.

\begin{tcolorbox}[colback=yellow!5!white,colframe=orange!80!black,
  title=Sufficient Decrease Lemma (Stated; Proof in Next Lecture)]
Suppose $f\in\mathcal{C}^{1,1}_L(\R^n)$ and let $\{x_k\}_{k\geq 0}$ be
generated by gradient descent with one of:
\begin{itemize}
  \item constant step size $t_k=\tfrac{1}{L}$;
  \item exact line search;
  \item backtracking with parameters $s\in\R_{++}$, $\alpha\in(0,\tfrac{1}{2})$, $\beta\in(0,1)$.
\end{itemize}
Then for any $x\in\R^n$ and $t>0$:
\[
  f(x) - f\!\bigl(x-t\grad f(x)\bigr) \geq M\norm{\grad f(x)}^2,
\]
where $M>0$ depends on the step-size strategy.
This guarantees that at every iteration the function value decreases
by at least a constant multiple of the squared gradient norm.
\end{tcolorbox}

%%% ----------------------------------------------------------------
\section{Summary: How the Results Connect}
%%% ----------------------------------------------------------------

\begin{tcolorbox}[colback=green!5!white,colframe=green!60!black,
  title=Lecture Flow and Key Connections]
\begin{enumerate}
  \item \textbf{Examples (Sec.~1)} empirically show that convergence varies with
    problem geometry $\to$ motivates a theoretical measure.
  \item \textbf{Quadratic convergence recurrence (Sec.~2)} reduces convergence
    to bounding the inner-product ratio $\rho_k$.
  \item \textbf{Kantorovich inequality (Sec.~3)} provides the universal lower bound
    $\rho_k\geq 4Mm/(M+m)^2$.
  \item \textbf{Linear convergence bound (Sec.~3)} follows immediately, with
    contraction factor $c=\bigl(\frac{M-m}{M+m}\bigr)^2$.
  \item \textbf{Condition number (Sec.~4)} extracts $\kappa=M/m$ as the key
    determinant; motivates feature normalisation and regularisation in ML.
  \item \textbf{Rosenbrock (Sec.~5)} demonstrates catastrophic ill-conditioning
    ($\kappa\approx2508$) on a simple 2D problem; establishes why it is
    a standard benchmark.
  \item \textbf{$L$-smoothness (Sec.~6)} generalises bounded curvature from
    quadratics (constant Hessian) to general smooth functions.
  \item \textbf{Descent Lemma (Sec.~7)} provides the quadratic upper envelope
    for $L$-smooth functions; it will be used next lecture to prove convergence
    of gradient descent for general smooth objectives.
\end{enumerate}
\end{tcolorbox}

%%% ----------------------------------------------------------------
\appendix
\section{Annexure: Additional Notes on Condition Number}\label{app:cond_annexure}

\noindent\textbf{Condition number of sublevel sets (Boyd \& Vandenberghe~\cite{boyd}, \S9.1).}
A more general geometric view defines the condition number in terms of the
\emph{width} of a convex set $C\subseteq\R^n$ in direction $q$ ($\norm{q}=1$):
\[
  W(C,\,q) = \sup_{z\in C}\,q^\top z \;-\; \inf_{z\in C}\,q^\top z.
\]
The \textbf{condition number of $C$} is then
\[
  \mathrm{cond}(C) = \frac{W_{\max}^2}{W_{\min}^2},
  \qquad
  W_{\max} = \sup_{\norm{q}=1} W(C,q),\quad
  W_{\min} = \inf_{\norm{q}=1} W(C,q).
\]
For the ellipsoid $\mathcal{E} = \{x \mid (x-x_0)^\top A^{-1}(x-x_0)\leq 1\}$
with $A\succ 0$, the width in direction $q$ is $2\norm{A^{1/2}q}$, giving
$W_{\min}=2\lambda_{\min}(A)^{1/2}$ and $W_{\max}=2\lambda_{\max}(A)^{1/2}$, so:
\[
  \mathrm{cond}(\mathcal{E}) = \frac{\lambda_{\max}(A)}{\lambda_{\min}(A)} = \kappa(A).
\]
Thus the condition number of an ellipsoidal sublevel set \emph{equals} the condition
number of the matrix defining it. For a general strongly convex $f$ satisfying
$mI \preceq \nabla^2 f(x) \preceq MI$, the $\alpha$-sublevel set
$C_\alpha = \{x \mid f(x)\leq\alpha\}$ is sandwiched between two balls
with radii proportional to $1/\sqrt{M}$ and $1/\sqrt{m}$ respectively,
giving the bound $\mathrm{cond}(C_\alpha)\leq M/m = \kappa$.
A small condition number means $C_\alpha$ is nearly spherical (all widths similar);
a large condition number means $C_\alpha$ is elongated, exactly the geometry that
slows gradient descent.

Figure~\ref{fig:boyd_zigzag} (reproduced from Boyd \& Vandenberghe~\cite{boyd},
Figure~9.2) shows gradient descent with exact line search on
$f(x)=\tfrac{1}{2}(x_1^2+10\,x_2^2)$, whose sublevel sets are ellipses with
$\kappa=10$. Even with this moderate condition number the iterates visibly
zigzag, bouncing between the steep and shallow directions of the ellipsoid.

\begin{figure}[h]
  \centering
  \includegraphics[width=0.52\textwidth]{boyd_fig9_2_zigzag.png}
  \caption{GD with exact line search on $f(x)=\tfrac{1}{2}(x_1^2+10x_2^2)$,
    $\kappa=10$ (Boyd \& Vandenberghe~\cite{boyd}, Fig.~9.2).
    The elliptical contours ($\kappa=10$) force the iterates into a characteristic
    zigzag pattern: each gradient step overshoots along the narrow direction
    and undershoots along the wide direction, slowing convergence markedly
    compared to the $\kappa=1$ (circular) case.}
  \label{fig:boyd_zigzag}
\end{figure}

\subsection{Effect on Convergence Speed}

If $\kappa\gg 1$:
\[
  c = \left(\frac{\kappa-1}{\kappa+1}\right)^2 \approx \left(1 - \frac{2}{\kappa}\right)^2 \approx 1,
\]
so $f(x_{k+1})\lesssim f(x_k)$: each step gives negligible improvement.
If $\kappa\approx 1$: $c\approx 0$ and convergence is near-instantaneous.

\begin{example}[Revisiting Examples~\ref{ex:gd_quad1} and~\ref{ex:gd_quad2}]
For $f=x^2+2y^2$: Hessian $A=\operatorname{diag}(2,4)$,
$\kappa=4/2=2$, $c=(1/3)^2\approx 0.11$ --- fast convergence (2--13 steps).

\noindent For $f=x^2+\tfrac{1}{100}y^2$: Hessian $A=\operatorname{diag}(2,\tfrac{1}{50})$,
$\kappa=2/(1/50)=100$, $c=(99/101)^2\approx 0.961$ --- slow convergence (201 steps).
\end{example}

Figure~\ref{fig:cond_comparison} shows convergence curves for several values of $\kappa$
and the theoretical contraction factor $c$ as a function of $\kappa$.

\begin{figure}[h]
  \centering
  \includegraphics[width=0.95\textwidth]{condition_number_comparison.pdf}
  \caption{Left: convergence of gradient descent on $f(x,y)=x^2+\kappa^{-1}y^2$
    for different condition numbers. Right: the theoretical contraction factor
    $c=\bigl(\frac{\kappa-1}{\kappa+1}\bigr)^2$ as a function of $\kappa$; as $\kappa$
    grows, $c\to 1$ and convergence stalls.}
  \label{fig:cond_comparison}
\end{figure}

\subsection{Role of Normalisation in Machine Learning}

\begin{tcolorbox}[colback=blue!5!white,colframe=blue!60!black,
  title=Practical Remark: Feature Normalisation and Condition Number]
In supervised learning the loss Hessian involves $H=X^\top X$
(where $X$ is the design matrix). If features have vastly different scales
(e.g.\ income in thousands vs.\ binary flags), the eigenvalues of $H$ can
span many orders of magnitude, giving $\kappa(H)\gg 1$ and very slow gradient descent.

\medskip
\textbf{Normalising} (mean-centering and scaling each feature to unit variance)
brings the eigenvalues of $H$ closer together, drastically reducing $\kappa(H)$
and accelerating convergence.

\medskip
\textbf{Degenerate case.} If $\lambda_{\min}(H)\approx 0$ (nearly collinear
features or underdetermined systems), the exact line-search step
$t_k = d_k^\top d_k/(2d_k^\top Hd_k)$ blows up along near-null eigenvectors,
causing highly unstable behaviour. $L_2$-regularisation (ridge) replaces $H$
by $H+\lambda I$, bounding $\kappa$ away from infinity at the cost of
a small bias.
\end{tcolorbox}

\section{Python Code for All Figures}
%%% ----------------------------------------------------------------

\begin{Verbatim}[fontsize=\small, label=Listing 1: Rosenbrock contour + GD iterates,
                  frame=single, labelposition=topline]
import numpy as np
import matplotlib; matplotlib.use('Agg')
import matplotlib.pyplot as plt

def rosenbrock(x):
    return 100*(x[1]-x[0]**2)**2 + (1-x[0])**2

def grad_rosenbrock(x):
    return np.array([
        -400*x[0]*(x[1]-x[0]**2) - 2*(1-x[0]),
         200*(x[1]-x[0]**2)
    ])

def backtrack(f, g, x, s=2, alpha=0.25, beta=0.5):
    t = s
    for _ in range(100):
        if f(x - t*g(x)) <= f(x) - alpha*t*(g(x)@g(x)):
            break
        t *= beta
    return t

def gradient_descent(f, g, x0, max_iter=10000, tol=1e-5):
    x = x0.copy(); path = [x.copy()]
    for _ in range(max_iter):
        grad = g(x)
        if np.linalg.norm(grad) < tol: break
        t = backtrack(f, g, x)
        x = x - t*grad; path.append(x.copy())
    return x, path

x0 = np.array([2.0, 5.0])
_, path = gradient_descent(rosenbrock, grad_rosenbrock, x0)
path = np.array(path)

x1, x2 = np.linspace(-3,3,500), np.linspace(-1,6,500)
X1, X2 = np.meshgrid(x1, x2)
Z = 100*(X2-X1**2)**2 + (1-X1)**2

fig, ax = plt.subplots(figsize=(7, 5))
ax.contour(X1, X2, Z, levels=np.logspace(0,3.5,30), cmap='viridis')
ax.plot(path[:,0], path[:,1], 'r-', lw=0.5, alpha=0.6, label='GD iterates')
ax.plot(1, 1, 'k*', ms=12, label='Optimum (1,1)')
ax.plot(x0[0], x0[1], 'go', ms=8, label='Start (2,5)')
ax.set_xlabel(r'$x_1$'); ax.set_ylabel(r'$x_2$')
ax.set_title('Rosenbrock: Contour Lines and GD Path')
ax.legend(); plt.tight_layout()
plt.savefig('rosenbrock_contour.pdf', dpi=150); plt.close()
\end{Verbatim}

\begin{Verbatim}[fontsize=\small, label=Listing 2: Condition number effect on convergence,
                  frame=single, labelposition=topline]
import numpy as np
import matplotlib; matplotlib.use('Agg')
import matplotlib.pyplot as plt

def gd_quadratic(kappa, max_iter=500, tol=1e-10):
    A = np.diag([1.0, 1.0/kappa]); x = np.array([1.0, 1.0])
    f_vals = [float(x@A@x)]
    for _ in range(max_iter):
        g = 2*A@x; d = -g
        t = (d@d) / (2*(d@A@d)); x = x + t*d
        fk = float(x@A@x); f_vals.append(fk)
        if fk < tol: break
    return f_vals

fig, axes = plt.subplots(1, 2, figsize=(12, 4))
for kappa, col in zip([1.5, 5, 20, 100],
                       plt.cm.plasma(np.linspace(0.1, 0.9, 4))):
    axes[0].semilogy(gd_quadratic(kappa), color=col,
                     label=rf'$\kappa={kappa}$', lw=2)
axes[0].set_xlabel('Iteration $k$')
axes[0].set_ylabel(r'$f(x_k)$ (log scale)')
axes[0].set_title('Convergence vs. Condition Number')
axes[0].legend(); axes[0].grid(True, alpha=0.3)

kap = np.linspace(1.01, 150, 400)
c   = ((kap-1)/(kap+1))**2
axes[1].plot(kap, c, 'b-', lw=2)
axes[1].axhline(1, color='r', ls='--', lw=1, label='$c=1$ (no progress)')
axes[1].set_xlabel(r'Condition Number $\kappa$')
axes[1].set_ylabel(r'Contraction factor $c$')
axes[1].set_title('Contraction Factor vs. Condition Number')
axes[1].legend(); axes[1].grid(True, alpha=0.3)
plt.tight_layout()
plt.savefig('condition_number_comparison.pdf', dpi=150); plt.close()
\end{Verbatim}

\begin{Verbatim}[fontsize=\small, label=Listing 3: Descent Lemma illustration,
                  frame=single, labelposition=topline]
import numpy as np
import matplotlib; matplotlib.use('Agg')
import matplotlib.pyplot as plt

L  = 3.0; x0 = 1.5
f  = lambda x: np.sin(x) + 0.5*x**2
df = lambda x: np.cos(x) + x
xr = np.linspace(-0.5, 3.5, 500)
Qp = f(x0) + df(x0)*(xr-x0) + 0.5*L*(xr-x0)**2
Qm = f(x0) + df(x0)*(xr-x0) - 0.5*L*(xr-x0)**2

fig, ax = plt.subplots(figsize=(7, 4))
ax.plot(xr, f(xr), 'b-', lw=2.5, label=r'$f(y)$')
ax.plot(xr, Qp,    'r-', lw=2,   label=r'$Q^+(y)$ upper bound')
ax.plot(xr, Qm,    'g--',lw=2,   label=r'$Q^-(y)$ lower bound')
ax.axvline(x0, color='gray', ls=':', lw=1.5)
ax.scatter([x0], [f(x0)], color='k', zorder=5, s=60,
           label=r'Expansion point $x$')
ax.set_ylim(-3, 9); ax.set_xlabel('$y$')
ax.set_title(r'Descent Lemma: $Q^-(y)\leq f(y)\leq Q^+(y)$')
ax.legend(fontsize=9); ax.grid(True, alpha=0.3)
plt.tight_layout()
plt.savefig('descent_lemma.pdf', dpi=150); plt.close()
\end{Verbatim}

\begin{Verbatim}[fontsize=\small, label=Listing 4: Lipschitz vs L-smooth comparison,
                  frame=single, labelposition=topline]
import numpy as np
import matplotlib; matplotlib.use('Agg')
import matplotlib.pyplot as plt

x = np.linspace(-2, 2, 1000); eps = 0.1
f_abs  = np.abs(x)
f_sm   = np.sqrt(x**2 + eps)
df_sm  = x / np.sqrt(x**2 + eps)
sign_x = np.where(x>0, 1., np.where(x<0, -1., 0.))

fig, axes = plt.subplots(1, 2, figsize=(11, 4))
axes[0].plot(x, f_abs, 'b-', lw=2.5,
             label=r'$|x|$ (Lipschitz, not $L$-smooth)')
axes[0].plot(x, f_sm,  'r-', lw=2,
             label=r'$\sqrt{x^2+0.1}$ ($L$-smooth, $L=1$)')
axes[0].set_title('Function Values')
axes[0].set_xlabel('$x$'); axes[0].legend(fontsize=9)
axes[0].grid(True, alpha=0.3)

axes[1].plot(x, sign_x, 'b-', lw=2.5,
             label=r'$\partial|x|$ (discontinuous at 0)')
axes[1].plot(x, df_sm,  'r-', lw=2,
             label=r'$\nabla\sqrt{x^2+0.1}$ (Lipschitz, continuous)')
axes[1].set_title('Gradient / Subgradient Comparison')
axes[1].set_xlabel('$x$'); axes[1].legend(fontsize=9)
axes[1].grid(True, alpha=0.3)
plt.tight_layout()
plt.savefig('lipschitz_vs_smooth.pdf', dpi=150); plt.close()
\end{Verbatim}

\begin{thebibliography}{9}
\bibitem{beck}
  A.\ Beck, \textit{Introduction to Nonlinear Optimization: Theory, Algorithms,
  and Applications with MATLAB}. SIAM, 2014.
\bibitem{bertsekas}
  D.\ Bertsimas and J.N.\ Tsitsiklis, \textit{Introduction to Linear Optimization}.
  Athena Scientific, 1997.
\bibitem{boyd}
  S.\ Boyd and L.\ Vandenberghe, \textit{Convex Optimization}.
  Cambridge University Press, 2004. pp.\ 461--475.
  Available: \url{https://web.stanford.edu/~boyd/cvxbook/}
\end{thebibliography}

\end{document}